\documentclass[10pt,a4paper]{amsart}
\usepackage[margin=19mm]{geometry}
\usepackage{amsmath,amssymb,mathtools}
\usepackage[scr=rsfs,cal=euler]{mathalfa}
\usepackage{xfrac}
\usepackage[unicode,hidelinks,bookmarks=false]{hyperref}
\newcommand{\ie}{i.e.\ }
\newcommand{\cf}{cf.\ }
\newcommand{\resp}{respectively\ }
\newcommand{\Subsection}{\subsection}
\providecommand{\nobreakdash}{\nobreak\hbox{-}}
\setlength{\emergencystretch}{2em}
\multlinegap0pt
\usepackage{bm}
\usepackage{bbm}
\usepackage{dsfont}
\usepackage{xspace}
\usepackage{cancel}
\usepackage{mathtools}
\usepackage{amsthm}
\makeatletter
\@ifundefined{thm}{\newtheorem{thm}{Thm}}{}
\makeatother
\title[The Corona Problem for Slice Holomorphic Functions via the M]{The Corona Problem for Slice Holomorphic Functions via the Matrix Corona Problem}
\author{Fabrizio Colombo, Elodie Pozzi, Irene Sabadini, Brett D. Wick}
\date{Source: arXiv:2608.12258v1 (CC BY 4.0)}
\begin{document}
\maketitle

\noindent\emph{Source and licence.} The Corona Problem for Slice Holomorphic Functions via the Matrix Corona Problem. Version arXiv:2608.12258v1 (CC BY 4.0), distributed under the Creative Commons Attribution 4.0 International licence, as recorded in the arXiv licence metadata for this exact version and shown on the abstract page https://arxiv.org/abs/2608.12258v1. Full rights notice: licenses/arxiv-2608.12258-CC-BY-4.0.txt.\par\smallskip

\noindent\textit{Editorial scope.} Counted unit: the Thm whose source label is \texttt{thm:detN-CB} in the article (source0.tex lines 2157--2177, source label \texttt{thm:detN-CB}), delivered together with the article's own complete original proof of it (source0.tex lines 2179--2261, one \texttt{proof} environment of 83 lines). No mathematics is written, repaired or re-derived: statement and proof are the article's own text line for line. Required context retained uncounted: none. Every definition, display and label of those lines is retained unchanged. Omitted from this package and not redistributed: 1--2156, 2178--2178, 2262--3056. Nothing inside a retained excerpt is omitted or altered: every retained body is delivered exactly as the article's lines are written, so no omission receipt of any kind accompanies this package. Citations: the retained excerpts carry no citation command at all, so no bibliography entry of the article is transcribed and no reference is redistributed. Reference adaptation: none was needed and none was made; every reference inside the delivered unit resolves inside it, so no unresolved reference or citation is produced. Printed slips kept literal: no printed word, symbol, hypothesis or number of the counted statement or of its proof was altered. Layout and class: the article's own class is \texttt{amsart} with packages including amsmath, amssymb, amsthm, latexsym, amsfonts, xcolor, hyperref, amsrefs, mathrsfs, multirow; the wrapper is the lane's pinned \texttt{amsart} editorial wrapper at 10pt on A4, which restates the theorem-like environments the retained text needs and adds the article's own macro definitions that the retained text needs (none), with extra wrapper packages (bm, bbm, dsfont, xspace, cancel, mathtools). The delivered numbering is the unit's own. Assets: no retained span contains a figure environment, a graphics-inclusion command, a PDF-page inclusion, a table or a TikZ picture; no image file is copied into this package, the package declares no asset dependency and no image file is redistributed with it. Licence: version arXiv:2608.12258v1 of this article is distributed under the Creative Commons Attribution 4.0 International licence, as recorded in the arXiv licence metadata for this exact version and shown on the abstract page https://arxiv.org/abs/2608.12258v1. A full rights notice is delivered at licenses/arxiv-2608.12258-CC-BY-4.0.txt. Characters: every retained excerpt is pure ASCII, so the delivered engine, XeLaTeX with its default Latin Modern text font, reports no missing character.
\begin{thm}\label{thm:detN-CB}
Let $F_1,\ldots,F_N$ and $G_1,\ldots,G_N$ be holomorphic functions on $\mathbb{D}$, and write $\hat{h}(z) := \overline{h(\overline{z})}$ for the Schwarz reflection (so $\hat h$ is holomorphic). Define the $2 \times 2N$ matrix
\begin{equation}\label{eq:tildeF-N}
\mathsf{F}(z) :=
\begin{pmatrix}
\begin{array}{cccccc}
F_1(z) & -G_1(z) & \cdots & F_N(z) & -G_N(z) \\
\hat G_1(z) & \hat F_1(z) & \cdots & \hat G_N(z) & \hat F_N(z)
\end{array}
\end{pmatrix},
\end{equation}
and let $\mathcal{M}(z) := \mathsf{F}(z)\mathsf{F}^*(z)$ be the associated $2 \times 2$ matrix. Then
\begin{equation}\label{eq:detN-CB}
\begin{aligned}
\det\mathcal{M}(z)
&= \sum_{r=1}^{N}\sum_{j=1}^{N} \bigl| F_j(z)\hat F_r(z) + G_r(z)\hat G_j(z)\bigr|^2 \\
&\quad + \sum_{1 \le r < j \le N} \bigl| F_j(z)\hat G_r(z) - F_r(z)\hat G_j(z)\bigr|^2 \\
&\quad + \sum_{1 \le r < j \le N} \bigl| G_j(z)\hat F_r(z) - G_r(z)\hat F_j(z)\bigr|^2.
\end{aligned}
\end{equation}
\end{thm}

\begin{proof}
Throughout the proof we suppress the variable $z$. The Gram matrix is
\[
\mathcal{M} =  \mathsf{F}\mathsf{F}^* =
\begin{pmatrix}
\displaystyle\sum_{j=1}^{N}\bigl(|F_j|^2 + |G_j|^2\bigr) & \displaystyle\sum_{j=1}^{N}\bigl(F_j\,\overline{\hat G_j} - G_j\,\overline{\hat F_j}\bigr)
 \\[0.6em]
\displaystyle\sum_{j=1}^{N}\bigl(\overline{F_j}\,\hat G_j - \overline{G_j}\,\hat F_j\bigr) &
\displaystyle\sum_{j=1}^{N}\bigl(|\hat F_j|^2 + |\hat G_j|^2\bigr)
\end{pmatrix},
\]
which matches the $2 \times 2$ matrix whose determinant we wish to compute (after using $|\hat F_j(z)|^2 = |F_j(\overline z)|^2$ and $\overline{\hat G_j(z)} = G_j(\overline z)$).

\bigskip

The key tool that will be used is the Cauchy--Binet formula, which we recall.  For any $m \times n$ matrix $A$ with $m \le n$,
\begin{equation}\label{eq:CB}
\det(A A^*) = \sum_{\substack{S \subseteq \{1,\ldots,n\} \\ |S| = m}} \bigl|\det A_S\bigr|^2,
\end{equation}
where $A_S$ is the $m \times m$ submatrix obtained by selecting the rows of $A^*$ indexed by $S$, equivalently, the columns of $A$ indexed by $S$ when $A$ is read row-wise. Applied to $\mathsf{F}$ (which is $2 \times 2N$), with $m = 2$, this gives
\begin{equation}\label{eq:CB-applied}
\det\mathcal{M} = \det( \mathsf{F} \mathsf{F}^*) = \sum_{1 \le i < k \le 2N} \bigl|\det F_{(i,k)}\bigr|^2,
\end{equation}
where $ F_{(i,k)}$ denotes the $2 \times 2$ submatrix formed by columns $i$ and $k$ of $F$.

Observe that the total number of such column pairs is $\binom{2N}{2} = 2N^2 - N$.

\bigskip

To facilitate the computation we split the corresponding $2 \times 2$ minors into certain collections. Index the columns of $\mathsf{F}$ as
\begin{align*}
\text{(odd)} \quad 1,\ldots,N: & \ \text{columns of the form } (F_j, \hat G_j)^{\top},\\
\text{(even)} \quad 1,\ldots,N:&  \ \text{columns of the form } (-G_{j}, \hat F_{j})^{\top}.
\end{align*}
A column pair $(i,k)$ falls into one of three cases.

\smallskip

\textit{Case 1: Both columns from the odd collection.} For $1 \le r < j \le N$,
\[
\det \begin{pmatrix} F_r & F_j  \\ \hat G_r & \hat G_j \end{pmatrix}
= F_r \hat G_j - F_j \hat G_r.
\]
Up to sign this is $F_j \hat G_r - F_r \hat G_j$, and the squared modulus is invariant. The contribution from these $\binom{N}{2}$ minors is
\[
\sum_{1 \le r < j \le N} |F_j \hat G_r - F_r \hat G_j|^2.
\]

\smallskip

\textit{Case 2: Both columns from the even collection.} For $1 \le r < j \le N$,
\[
\det \begin{pmatrix} -G_r & -G_j \\ \hat F_r  & \hat F_j \end{pmatrix}
= -G_r \hat F_j + G_j \hat F_r = G_j \hat F_r - G_r \hat F_j.
\]
The contribution from these $\binom{N}{2}$ minors is
\[
\sum_{1 \le r < j \le N} |G_j \hat F_r - G_r \hat F_j|^2.
\]

\smallskip

\textit{Case 3: One column from each collection.} For $r \in \{1,\ldots,N\}$ (odd) and $j \in \{1,\ldots,N\}$ (even),
\[
\det \begin{pmatrix} F_r & -G_j \\ \hat G_r & \hat F_j \end{pmatrix}
= F_r \hat F_j + G_j \hat G_r.
\]
This gives $N^2$ minors (one per ordered pair $(r,j) \in \{1,\ldots,N\}^2$). Their contribution is
\[
\sum_{r=1}^{N}\sum_{j=1}^{N} |F_r \hat F_j + G_j \hat G_r|^2.
\]

\bigskip

This is assembled by substituting the three cases into \eqref{eq:CB-applied},
\[
\det\mathcal{M}
= \sum_{r,j=1}^{N} |F_r \hat F_j + G_j \hat G_r|^2
+ \sum_{1\leq r<j\leq N} |F_j \hat G_r - F_r \hat G_j|^2
+ \sum_{1\leq r<j\leq N} |G_j \hat F_r - G_r \hat F_j|^2,
\]
which agrees with \eqref{eq:detN-CB} after the trivial relabeling $(r,j) \leftrightarrow (j,r)$ in the first sum (which runs over all ordered pairs, so the sum is unchanged).  The minor counts are $\binom{N}{2} + \binom{N}{2} + N^2 = 2N^2 - N = \binom{2N}{2}$, as required.
\end{proof}

\end{document}
